% Options for packages loaded elsewhere
\PassOptionsToPackage{unicode}{hyperref}
\PassOptionsToPackage{hyphens}{url}
\documentclass[
  ignorenonframetext,
]{beamer}
\newif\ifbibliography
\usepackage{pgfpages}
\setbeamertemplate{caption}[numbered]
\setbeamertemplate{caption label separator}{: }
\setbeamercolor{caption name}{fg=normal text.fg}
\beamertemplatenavigationsymbolsempty
% remove section numbering
\setbeamertemplate{part page}{
  \centering
  \begin{beamercolorbox}[sep=16pt,center]{part title}
    \usebeamerfont{part title}\insertpart\par
  \end{beamercolorbox}
}
\setbeamertemplate{section page}{
  \centering
  \begin{beamercolorbox}[sep=12pt,center]{section title}
    \usebeamerfont{section title}\insertsection\par
  \end{beamercolorbox}
}
\setbeamertemplate{subsection page}{
  \centering
  \begin{beamercolorbox}[sep=8pt,center]{subsection title}
    \usebeamerfont{subsection title}\insertsubsection\par
  \end{beamercolorbox}
}
% Prevent slide breaks in the middle of a paragraph
\widowpenalties 1 10000
\raggedbottom
\AtBeginPart{
  \frame{\partpage}
}
\AtBeginSection{
  \ifbibliography
  \else
    \frame{\sectionpage}
  \fi
}
\AtBeginSubsection{
  \frame{\subsectionpage}
}
\usepackage{iftex}
\ifPDFTeX
  \usepackage[T1]{fontenc}
  \usepackage[utf8]{inputenc}
  \usepackage{textcomp} % provide euro and other symbols
\else % if luatex or xetex
  \usepackage{unicode-math} % this also loads fontspec
  \defaultfontfeatures{Scale=MatchLowercase}
  \defaultfontfeatures[\rmfamily]{Ligatures=TeX,Scale=1}
\fi
\usepackage{lmodern}
\ifPDFTeX\else
  % xetex/luatex font selection
\fi
% Use upquote if available, for straight quotes in verbatim environments
\IfFileExists{upquote.sty}{\usepackage{upquote}}{}
\IfFileExists{microtype.sty}{% use microtype if available
  \usepackage[]{microtype}
  \UseMicrotypeSet[protrusion]{basicmath} % disable protrusion for tt fonts
}{}
\makeatletter
\@ifundefined{KOMAClassName}{% if non-KOMA class
  \IfFileExists{parskip.sty}{%
    \usepackage{parskip}
  }{% else
    \setlength{\parindent}{0pt}
    \setlength{\parskip}{6pt plus 2pt minus 1pt}}
}{% if KOMA class
  \KOMAoptions{parskip=half}}
\makeatother
\setlength{\emergencystretch}{3em} % prevent overfull lines
\providecommand{\tightlist}{%
  \setlength{\itemsep}{0pt}\setlength{\parskip}{0pt}}
% Slide layout defaults for warning-clean scientific decks.
\usepackage{etoolbox}
\IfFileExists{xurl.sty}{\usepackage{xurl}}{}
\IfFileExists{seqsplit.sty}{\usepackage{seqsplit}}{\newcommand{\seqsplit}[1]{#1}}
\protected\def\breakseq#1{\seqsplit{#1}}
\protected\def\breaktt#1{\begingroup\ttfamily\seqsplit{#1}\endgroup}
\setlength{\emergencystretch}{6em}
\tolerance=5000
\hbadness=10000
\hfuzz=1pt
\setlength{\tabcolsep}{2pt}
\AtBeginEnvironment{longtable}{\tiny\renewcommand{\arraystretch}{0.86}\setlength{\tabcolsep}{1pt}}
\AtBeginEnvironment{tabular}{\tiny\renewcommand{\arraystretch}{0.86}\setlength{\tabcolsep}{1pt}}

% Natbib and cross-reference fallbacks — slides don't load natbib
% or cleveref, but combined-PDF manuscript prose may emit these
% commands. The fallback renders citations as a bracketed key list
% and cross-references as detokenized labels so slides don't fail on
% undefined control sequences or raw underscores. \providecommand is
% a no-op if packages load later.
\providecommand{\citep}[1]{[#1]}
\providecommand{\citet}[1]{#1}
\providecommand{\citealp}[1]{#1}
\providecommand{\citeauthor}[1]{#1}
\providecommand{\citeyear}[1]{#1}
\providecommand{\cref}[1]{\texttt{\detokenize{#1}}}
\providecommand{\Cref}[1]{\texttt{\detokenize{#1}}}

\newtheorem{proposition}{Proposition}
\newtheorem{hypothesis}{Hypothesis}
\usepackage{bookmark}
\IfFileExists{xurl.sty}{\usepackage{xurl}}{} % add URL line breaks if available
\urlstyle{same}
% fallback for those not using the hyperref driver hyperxmp:
\makeatletter
\@ifundefined{xmpquote}{\newcommand{\xmpquote}[1]{#1}}{}
\makeatother
\hypersetup{
  hidelinks,
  pdfcreator={LaTeX via pandoc}}

\author{\texorpdfstring{}{}}
\date{}

\begin{document}

\section{Abstract}\label{sec:abstract}

\begin{frame}[allowframebreaks]{Abstract}
DigiPPPiP extends Partner Pen Play in Parallel (PPPiP) from a co-present
paper practice into a reproducible framework for cyberphysical, remote,
semisynchronous, asynchronous, accessible, and place-responsive dyadic
drawing \breakseq{[@mikhailova2018pppip]}. The framework is positioned within
collaborative-work studies of shared drawing surfaces
\breakseq{[@tang1991collaborativework]}, cyber-physical-social systems
\breakseq{[@sobb2023cpss]}, mediated embodiment and presence
\breakseq{[@biocca1997cyborg; @lee2004presence]}, digital intimacy technologies
\breakseq{[@hassenzahl2012love; @neustaedter2012intimacy]}, and cautious
hyperscanning methodology \breakseq{[@czeszumski2020hyperscanning;
@hamilton2021hyperscanning]}.

The manuscript translates the project brief into a modular research
artifact: 9 temporal-spatial modalities, 13 evidence dimensions, 6
planned outcome measures, and 39 conceptual figures are generated or
registered by the project code rather than hand-maintained prose. The
computational layer is explicitly illustrative, not empirical: the
active-inference, inter-brain-synchrony, Forman-Ricci,
narrative-information, source-quality, accessibility, and neuroergonomic
outputs are deterministic conceptual models designed to make theoretical
commitments inspectable. Formal equations are consolidated in
\breakseq{[@sec:formalisms\_appendix]} so the main line can state evidence
boundaries before presenting mathematical detail.

The argument is that DigiPPPiP should be treated as human-human
relational technology. Its irreducible design kernel is modest: two
partners, a shared mark field, perceptible traces of agency, a temporal
relation among contributions, and consentful control over persistence.
Digital infrastructure expands the substrate of shared mark-making
without replacing the second-person, embodied, narrative, and affective
conditions that make the practice relationally meaningful
\breakseq{[@schilbach2013secondperson; @bolis2024secondperson;
@vaisvaser2024neurodynamics]}. The manuscript explicitly avoids claims
that DigiPPPiP is clinically therapeutic, universally accessible, or
causally validated by neural synchrony; those remain empirical questions
for controlled studies. Cross-references use Pandoc labels, so sections,
equations, figures, and tables remain automatically numbered across PDF
and web render targets.
\end{frame}

\end{document}
